\section{Dividing cycles, Yoccoz puzzle, and renormalization}\label{principal nest}

Let $(f: U'\ra U)$ be a quadratic-like map.
We assume that the domains $U'$ and $U$ are smooth disks,
$f$ is even,  
and we  normalize $f$ so that $0$ is its critical point.  
% Our standing assumption is that {\it the Julia set $J(f)$ is connected 
% and both fixed points of $f$ are repelling}.                      
We let $U^m = f^{-m}(U)$. The boundary of $U^m$ is called  
the {\it equipotential of depth $m$}.   

By means of straightening, we can define external rays for $f$. They form a foliation of $U\sm K(f)$
transversal   to the equipotential $\di U$. Each ray is labeled by its {\it external angle}. 
These rays will play purely combinatorial role, so particular choice of the straightening
is not important. 
  
\subsection{Dividing cycles and associated Yoccoz puzzles}\label{sec: first king} 

  Let us consider a repelling periodic point $\gamma$ of period $\bt$ and the correspoding cycle 
$\bgamma=\{f^k\gamma\}_{k=0}^\bt$. This point (and the cycle) is called {\it dividing}
if there exist at least two rays landing at it.  
For instance, the landing point of the zero ray is a non-dividing fixed point, 
while the other fixed point is dividing (if repelling). 

{\it In what follows, we assume that $\gamma$ is dividing.}
Let $\RR(\gamma)$ (resp., $\RR(\bgamma)$) stand for the family of rays landing at $\gamma$ (resp., $\bgamma$).
Let ${\bf s} =\# \RR(\gamma)$ and let $\r= \bt\s= \# \RR(\bgamma)$. 
These rays divide $U$ into $\bt(\s-1)+1$ closed topological disks $Y^0(j)\equiv Y^0_\sbg(j)$ called 
{\it Yoccoz puzzle pieces of depth $0$}.     


{\it Yoccoz puzzle pieces $Y^m(j)\equiv Y^m_\sbg(j)$ of depth $m$} are  defined as the 
pullbacks of $Y^0(i)$ under $f^m$. 
They tile the neighborhood of $K(f)$ bounded by the equipotential $\di U^m $.
Each of them is bounded by finitely many arcs of this equipotential and finitely many  
external rays of $f^{-m}(\RR(\bgamma))$. 
We will also use notation $ Y^m(z)$ for the puzzle piece $Y^m(j)$ containing $z$ in its interior. 
If $f^m(0)\not\in \bgamma$, then there is a well defined {\it critical} puzzle piece  $Y^m\equiv Y^m(0)$.
The critical puzzle pieces are nested around the origin:
$$
   Y^0\supset Y^1\supset Y^2 \dots \ni 0. 
$$
Notice that all $Y^m$, $m\geq 1$, are symmetric with respect to the origin.  


Let us take a closer look at some puzzle piece $Y=Y^m(i)$.
Different arcs of $\di Y$  meet at the {\it corners} of $Y$.
The corners where two external rays meet will be called {\it vertices} of $Y$;
they are $f^m$-preimages of $\bgamma$.
Let $K_Y= K(f)\cap Y$. It is a closed connected set that meets the boundary $\di Y$
at its  vertices. Moreover, the external rays meeting at a vertex $v\in \di Y$ chop off from
$K(f)$ a continuum $S_Y^v$, the component of $K(f)\sm \inter Y$ containing $v$. 

Let $\YY_\sbg$ stand for the family of all puzzle pieces $Y_\sbg^m(j)$. 

Let us finish with an obvious observation that will be constantly exploited:

\begin{lem}\label{Subset}
  If a puzzle piece $Y^n(z)$ of $\YY_\bgamma$ does not touch the cycle $\bgamma$
then $Y^n(z)\Subset Y^0(z)$.
\end{lem}

\subsection{Renormalization associated with a dividing cycle}\label{renorm-sec}

\begin{lem}[see \cite{Th,M-rays}]\label{cv pp}
  The puzzle piece $X^0 \equiv Y^0(f(0))$ of $\YY_\bgamma$ 
containing the critical value has only one vertex, 
and thus is bounded by only two external rays (and one equipotential).
\end{lem} 
  
In what follows, 
$\gamma$ will denote the point of the cycle $\bgamma$ such that $f(\gamma)$ is the vertex of $X^0$.
Notice that $f(0)\in \inter X^0$ for otherwise $f(0)=f(\gamma)$, which is impossible since $0$ is the
only preimage of $f(0)$. 

Since the critical puzzle piece $Y^1$ is the pullback of $X^0$ under $f$, it has two vertices, 
$\gamma$ and $\gamma'= -\gamma$, and is bounded by four rays,
two of them landing at $\gamma$ and two landing at $\gamma'$.%
\footnote{We will usually say that ``a puzzle piece is bounded by several external rays'' 
without mentioning equipotentials that also form part of its boundary.} 

\begin{lem}(see \cite{D})
  Let $X^\r \equiv Y^\r (j)\subset X^0$  be the puzzle piece attached to the boundary of $X^0$.
Then $f^\r : X^\r \ra X^0$ is a double branched covering.
\end{lem}
 
\begin{pf}
    Let $\CC^0$ be the union of the two rays that bound $X^0$,
and let $\CC^\r  $ be $\CC^0$ cut by the equipotential $\di U^\r$.
Let us orient $\CC^0$ and then induce the orientation to $\CC^\r $.
    Since $f^\r : \CC^\r \ra \CC^0$ is an orientation preserving homeomorphism, 
it maps $X^\r $ onto $X^0$. 

Since for $m=1,\dots, \r - 1$,  the arcs $f^m(\CC^\r)\subset \di (f^m X^\r )$ are disjoint  from
$\inter X^0$, the puzzle pieces $f^m X^\r $ are not contained in $X^0$.
Since they have a bigger depth than $X^0$, they are disjoint from $\inter X^0$.
It follows that all the puzzle pieces $f^m (X^\r )$, $m=0,1,\dots, \r -1$, 
have pairwise disjoint interiors. 
(Otherwise $f^m(X^\r)\supset f^n (X^\r)$ for some $ \r> m > n\geq 0$,
and applying $f^{\r-m}$, we would conclude that $X^0\supset f^{\r-n+m} (X^\r)$.)    

Moreover, the puzzle piece  $f^{\r-1}X^\r=Y^1$ is  critical
since $Y^1$ is the only pullback of $X^0$ under $f$. 
Hence the puzzle pieces $f^m(X^\r)$, $m=0,1,\dots, \r-2$, do not contain $0$. 
It follows that  $\deg(f^\r : X^\r \ra X^0)=2$.
\end{pf}

\makefigure{pullback}{}

\begin{cor}\label{renorm}
  If $f(0)\in \inter X^\r $ then the puzzle piece $Y^{\r+1}$ has four vertices,
and the map $f^\r : Y^{\r+1}\ra Y^1$ is a double branched covering.  
\end{cor}

   Let $\Theta(\bgamma)\subset \T$ be the set of external angles of the rays of $\RR(\bgamma)$.
There is a natural equivalence relation on $\Theta(\bgamma)\subset \T$:
two angles are equivalent if the corresponding rays land at the same periodic point. 
Let us consider the hyperbolic convex hulls of these equivalence classes 
(in the disk $\D$ viewed as the hyperbolic plane). 
The union of the boundaries of these convex hulls 
is a  finite lamination $\PP=\PP(\bgamma)$ in $\D$ which is also called the 
{\it periodic ray portrait}. One can characterize all possible ray portraits that appear in this way
(see \cite{M-rays}). 


\begin{definition} 
A map $f$ is called $\PP$-{\it renormalizable} 
(or, ``renormalizable with combinatorics $\PP$'')
if $f(0)\in \inter X^\r$ and $f^{\r m}(0)\in  Y^{\r+1}$  
for all $m=0,1,2,\dots$. 
In this case, the double covering $f^\r : Y^{\r+1}\ra Y^1$ is called
the {\it renormalization} $R_\PP f = R_\sbg f$ of $f$ (associated with the cycle $\bgamma$).
The corresponding {\it  little (filled) Julia set} $\KK = K(R f)$ 
is defined as 
$$
        \{ z: f^{\r m}z\in Y^{\r + 1} ,\quad  m=0,1,2\dots\}
$$
If the little Julia sets $f^m \KK$, $k=0,1\dots, \r-1$, 
are pairwise disjoint, then the renormalization is called {\it primitive};
otherwise it is called {\it satellite}. 

In case $\bgamma$ is the dividing fixed point of $f$,
the map  is also called {\it immediately renormalizable}.
(This is a particular case of the satellite renormalization.)
\end{definition}



\begin{rem} The above  definition of renormalization is not quite standard since the map 
$R_\sbg f$ is not quadratic-like. To obtain the usual notion of renormalization, 
one should thicken the domain of $R_\sbg f$  a bit to make it quadratic-like (see \cite{D,M}). 
This thickenning does not change the Julia set, so $K(Rf)$ possesses all the properties of 
quadratic-like Julia sets. In particular, it has two fixed points, one of which is either
non-repelling or dividing.

Note also that in the case when $f^r(0)=\gamma'$ the puzzle piece
$Y^{\r+1}$ degenerates (is pinched at $0$), but this does not effect any of 
further considerations. 
\end{rem}

Given a periodic ray portrait $\PP$, 
the set of parameters $c\in \C$ for which the quadratic polynomial $P_c$ is $\PP$-renormalizable form a 
{\it little copy} $M_\PP$ of the Mandelbrot set (``$M$-copy''). Thus, there is one-to-one correspondence
between the admissible ray portraits and the little M-copies.
So, one can encode the combinatorics of the renormalization by the little $M$-copies themselves. 

\subsection{Molecule Condition}

The  molecule $\MM$ defined in the Introduction consists of the quadratic maps which are:

\nin\ssk $\bullet$
 either finitely many times renormalizable, all these renormalizations are satellite,
  and the last renormalization has a non-repelling cycle;

\nin\ssk $\bullet$
  or  infinitely many times renormalizable, with all the renormalizations satellite.

\ssk The molecule condition that we are about to introduce will  ensure that our map $f$
  has frequent ``qualified'' primitive renormalizations.
Though  $f$ is allowed to be satellite renormalizable once in a while,
{\it we will   record only the primitive renormalizations}.
They  are naturlly ordered according to their periods,
$1=p_0 <p_1<p_2<\dots$, where  $p_i$ is a multiple of $p_{i-1}$. 
% Moreover, the corresponding $M$-copies $M_i$ are nested:
% $M=M_0\supset M_1\supset M_2\supset\dots$ 

Along with these ``absolute'' periods of the primitive renormalizations,
we will consider {\it relative periods} $\tl p_i = p_i/p_{i-1}$ and the corresponding $M$-copies
$\tl M_i$ that encode the combinatorics of $R^i f$ as the renormalization of $R^{i-1} f$. 
% Any sequence of relative $M$-copies $M_i$ can appear as such combinatorics. 

% The renormalization level $i$ is called {\it satellite} if $R^{i-1} f$ is immediately renormalizable;
% otherwise it is called primitive. 

Given an $\eta>0$, we say that a sequence of primitively renormalizable quadratic-like maps $f_i$
satisfies the $\eta$-{\it molecule condition}
if the corresponding  $M$-copies $M_i$  stay $\eta$-away from the molecule $\MM$ 
(the latter is defined in the Introduction).
We say that $\{f_i\}$ satisfies the {\it molecule condition} if it does it for some $\eta>0$.
    
An infinitely primitively  renormalizable map $f$  satisfies the $\eta$-molecule condition
if the sequence of its primitive renormalizations $R^i f$ does
(i.e.,  the corresponding relative copies $\tl M_i$ stay $\eta$-away from the molecule $\MM$).   
(And similarly, for  the non-quantified molecule condition.)

\msk There is, however, a more specific combinatorial way to describe the molecule condition. 

Let us consider a quadratic-like map $f$ with straightening $P_c$, $c\in M$.
We are going to associate  to  $f$ (in some combinatorial region, and with some choice involved)
three combinatorial parameters, $(\r,\q,\n)$ (``period, valence, and escaping time'')
whose boundedness will be equivalent to the molecule condition.
 
\ssk
{\it Assume first that $f$ admits a dividing cycle $\bgamma$ with the ray portrait $\PP$
with $\r$ rays}. 
This happens if and only if $c$ belongs to the parabolic limb of $M$ cut off by
two external rays landing at an appropriate parabolic point. 
%(where an attracting cycle of period
%$\bt$ bifurcates into an attracting cycle of period $\bt\s=\r$).  
%However, unlike $f_*$, {\it the map  $f$ is not necessarily $\PP(\bgamma$)-renormalizable}. 

\ssk
On the central domain of $\C\sm (\RR(\bgamma)\cup \RR(\gamma'))$,
$f^\r$ has a unique fixed poin $\alpha$. 
{\it Next, we assume that $\alpha$ is repelling and there are $\q$ rays landing at it.} 

\ssk {\it Assume next that the finite orbit $f^{\r j(0)}$, $j=1,\dots, \q\n-1$,
      does not escape the central domain of  $\C\sm (\RR(\gamma)\cup \RR(\gamma'))$}.
       This happens if and only if 
  $c$ lies outside certain {\it decorations} (see \cite{decorations})
  of the above parabolic limb. In particular, this happens if $f$ is
  {\it satellite  $\PP(\gamma$)-renormalizable}.   

\ssk {\it Finally, assume that  $\n$ is the first moment $n$ such that  $f^{\r\q n}(0)$ escapes the
  central domain of $\C\sm (\RR(\alpha)\cup \RR(\alpha'))$}. 
   This happens if and only if 
  $c$ belongs to the union  of $2^\n$ {\it decorations} 
  inside the above parabolic limb. In particular, the map $f$ 
  {\it is not  $\PP(\balpha$)-renormalizable}.   

\ssk 
 Under the above assumptions, 
%  we have well defined parameters $(\r,\q,\n)$ associated with $f$.
  we say that $f$ satisfies the $(\bar\r, \bar\q, \bar\n)$-{\it molecule condition} if
  there is a choice of $(\r, \q,\n)$ with $\r\leq \bar\r$, $\q\leq \bar \q$ and $\n\leq \bar \n$.  


\begin{lem}
     The $(\r,\q,\n)$-molecule condition is equivalent to the $\eta$-molecule condition.%
\footnote{In the sense that if $f$ satisfies $(\r,\q,\n)$-molecule condition then it satisfies
 $\eta$-molecule condition with some $\eta=\eta(\r,\q,\n)$, and the other way around.}  
\end{lem}

\begin{pf}
  If $f$ satisfies $(\r,\q,\n)$-molecule condition then $c$ belongs to the finite union of decorations.
Each of them does not intersect the molecule $\MM$, so $c$ stays some distance $\eta$ away from $\MM$. 

Vice versa, assume there is a sequence of maps $f_i$ satisfying the $\eta$-molecule condition,
but with  $(\r,\q,\n)\to \infty$ for any choice of $(\r,\q,\n)$.
Let us select a convergent subsequence $c_i\to c$. Since $c\not\in \MM$, 
there can be only finitely many hyperbolic components $H_0,H_1,\dots,H_m$ of $\MM$
such that $H_0$ is bounded by the main cardioid, $H_{k+1}$ bifurcates from $H_k$,
and $f$ is $m$ times immediately  renormalizable with the corresponding combinatorics.   
Let us consider the two rays landing at the last bifurcation point (where $H_m$ is attached to $H_{m-1}$),
and the corresponding parabolic limb of the Mandelbrot set.%
\footnote{If $m=0$ then we consider the whole Mandelbrot set. } 
This parabolic point has certain period $\r$. 

Since the quadratic polynomial $P_c$ is not immediately renormalizable any more,
the corresponding cycle $\balpha$ is repelling with $\q$ rays landing at each of  its periodic points,
and there is some escaping time $\n$. 
So, $P_c$ satisfies $(\r,\q,\n)$-molecule condition.
Since this condition is stable under perturbations, $P_{c_i}$ satisfy it as well -- contradiction.  
\end{pf}

%{\it In what follows we assume that $\r\leq \bar\r$, $\q\leq \bar \q$ and $\n\leq \bar \n$ for
%some bounds $\bar\r$, $\bar\q$ and $\bar \n$ fixed throughout the paper.  
%All the constants will depend on these bounds, unless otherwise is explicitly stated.}

{\it In what follows we assume that parameters $\r,\q,\n$ are well defined for a map $f$ under consideration,}
so in particular, we have two dividing cycles, $\bgamma$ and $\balpha$. 
We let $\k=\r\q\n$.
Let us state for the record the following well-known combinatorial property:

\begin{lem}\label{zeta}
   The point $\zeta=f^\k(0)$ is separated from $\alpha$ and $0$ by the rays landing at $\alpha'$.%
\footnote{while the latter two points are not separated.}
\end{lem}

 \subsection{Combinatorial separation between $\gamma$ and $\balpha$}\label{comb sep} 
  
%  Assume now that $f$ is renormalizable with combinatoics $\PP$ associated with a dividing
%cycle $\bgamma$, but the renormalization $g =  R_\sbg f$ is not immediately renormalizable.
%Let $\alpha$ be the dividing fixed point of $g$, let $\balpha$ be its cycle (of period $\r$), 
%and let $\alpha'=-\alpha$.

Along with the puzzle $\YY_\sbg$ associated with $\bgamma$, let us consider the puzzle $\YY_\sba$
associated with $\balpha$. 
The critical puzzle piece $Y_\sba^1$ has two vertices, $\alpha$ and $\alpha'$,
and is bounded by four external rays landing at these vertices. Let $\CC$ 
be the union of the two external rays of $\di Y^1_\sba$ landing at $\alpha$, 
and let $\CC'$  be the symmetric pair of external rays landing at $\alpha'$. 
     
Recall that $\bt$ stands for the period of $\bgamma$ and $\s$ stands for $\# \RR(\bgamma)$. 

\begin{lem}\label{separation}
  There exist inverse branches $f^{-\bt m} |\, \CC'$, $m=0,1,\dots, \s-1$, 
such that the union of the arcs $f^{-\bt m} (\CC')$  separates $\gamma$ from 
the cycle $\balpha$ and the co-cycle $\balpha'$
(except that $\alpha'\in \CC'$). 
%Moreover, the puzzle piece $Y_\sba^{}(\gamma)$ has $\s$ vertices, and is bounded by
%the arcs  $f^{-\bt m} (\CC')$ (approperiately cut) and appropriate equipotentials.
\end{lem}

\begin{pf}
Let us pull the puzzle piece $Y_\sbg^1$ along the orbit $\bgamma$ 
(or equivalently, along the orbit  $\balpha$). 
By Corollary \ref{renorm}, 
the corresponding  inverse branches 
$$
   f^{-m} : Y^1_\sbg \ra f^{-m}(Y^1_\sbg) = f^{\r-m} Y^{\r+1}_\sbg,\quad  m=0,1,\dots, \r-1,
$$ 
have disjoint interiors.  
Hence each of these puzzle pieces contains exactly one point of $\balpha$. 
Moreover, non of these puzzle pieces except $Y^1_\sbg$ may intersect $\balpha'$
(for otherwise, its image would contain two points of $\balpha$).   

It follows from standard properties of quadratic maps 
that the arc $\CC'$ {\it separates} 
$\gamma$ (which is the non-dividing fixed point of $R_\sbg f$) from $\gamma'$ and
$\alpha$ (which is the dividing fixed point of $R_\sbg f$). 
Hence the arcs $f^{-\bt m} (\CC')$, 
separate $\gamma$ from  $f^{-\bt m}(\gamma')$ and $f^{-\bt m}(\alpha)$. 

Since each of the puzzle pieces $f^{-\bt m} ( Y^1_\sbg)$,  $m=0,1,\dots, \s-1$, has two vertices
($\gamma$ and $f^{-\bt m}(\gamma')$) and  their union forms a neighborhood of $\gamma$,
the rest of the Julia set is separated from $\gamma$ by the union of arcs
$f^{-\bt m}(\CC')$, $m=0,1,\dots, \s-1$. 
It follows that this union separates $\gamma$ from the whole cycle $\balpha$, 
and from the co-cycle $\balpha'$. 
\end{pf}

Together with Lemma \ref{Subset} this yields: 

\begin{cor} We have:
  $Y_\sba^{\r}(\gamma) \SUBset Y^0_\sba (\gamma)$ % \note{In fact, $\r-\bt+1$}
\end{cor}


\begin{pf}
   The puzzle piece $Y_\sba^{\r}(\gamma)$ contains $\gamma$ and does not cross the arcs 
$f^{-\bt m} (\CC')$, $m=0,1,\dots, \s-1$, from Lemma \ref{separation}.
Hence it does not intersect $\balpha$, and the the conclusion follows from Lemma \ref{Subset}.
\end{pf}


 \subsection{A non-degenerate annulus} 
%Let $\q = \# \RR(\alpha)$. 
%{\it Let us assume that the map $g=R_\sbg f$ is not immediately renormalizable.}
%Then there is an $\n$ such that $\zeta = g^{\q \n}(0)=f^{\r\q\n}(0)$ is separated from $\alpha$ by $\CC'$.
%Let $\r\q\n= \k$. 

 {\it In what follows we will be dealing only with the puzzle $\YY_\ba$, so we will skip the label $\ba$
in notation.} 
  
Since by Lemma \ref{zeta}, the point $\zeta =f^\k (0)$ is separated from $\alpha$ by $\CC'$,
the union of arcs $f^{-\bt m} (\CC')$, $m=0,1,\dots, \s-1$, from Lemma \ref{separation}
separates $\zeta$ from the whole cycle $\balpha$.
By Lemma \ref{Subset},  $Y^{\r}(\zeta)\Subset Y^0(\zeta)$.
Pulling this back by $ f^\k$,  we conclude:

\begin{lem}
We have: 
       $Y^{\k+\r} \Subset  Y^\k$.
\end{lem}

We let $E^0= Y^{\k+\r}$. 

\subsection{Buffers attached to the vertices of $P=Y^{\r\q (\n-1)+1}$}\label{P-sec}


 Let us consider a nest of critical puzzle pieces
$$
   Y^1 \supset  Y^{ \r\q +1 } \supset Y^{2\r\q+1 } \supset\dots \supset Y^{\r\q (\n-1)+1} = P. 
$$
 Since $f^{\r\q} : Y^{\r\q+1}\ra Y^1 $ is a double branched covering such that
$$
   f^{\r\q m}(0)\in Y^1, \quad m=0,1,\dots, \n-1,
$$
the puzzle piece $Y^{\r\q k+1}$  is mapped by 
 $f^{\r\q}$ onto $Y^{\r\q (k-1)+1}$ as a double branched covering, $k=1\dots,\n-1 $. 
However, since $f^\k(0)=f^{\r\q\n}(0)\not\in Y^1$, there are two non-critical puzzle pieces of depth 
$\k+1$ mapped univalently onto $P$ under $g^{\q}$.  
One of these puzzle pieces, called $Q_L$, is attached to the point $\alpha$,
another one, called $Q_R$, is attached to $\alpha'$.
The following Lemma is similar to Lemma  2.1 of \cite{decorations}:

\begin{lem}\label{Qv}
 For any vertex of $P$, there exists a puzzle piece $Q^v\subset P$ of
 depth $\r(2\n-1)\q + 1$ attached to the boundary rays of $P$ landing at $v$
 which is a univalent $f^\k$-pullback of $P$. Moreover, these puzzle pieces are
 pairwise disjoint.
\end{lem}


\subsection{Modified principle nest}\label{modified nest}
Until now, the combinatorics of the puzzle depended only on the parameters $(\r,\q,\n)$.
Now we will dive  into the deeper waters.

Let $l$ be the first return time of $0$ to $\inter E^0$ and let $E^1= Y^{\k+ \r +l}$
be the pullback of $E^0$ along the orbit $\{f^m (0)\}_{m=0}^l$.
Then $f^l: E^1\ra E^0$ is a double branched covering.

\begin{cor}\label{non-degenerate annulus}
     We have: $E^1\SUBset  E^0$.
\end{cor}

\begin{pf}
  Since $\{f^m(0)\}_{m=1}^{\r-1}$ is disjoint from $Y_\sbg^1\supset Y_\sba^1 \supset E^0$,
we have: $l\geq \r$. Hence 
$$
      f^l(E^0)\supset f^\r(E^0) \Supset E^0 = f^l(E^1),
$$
and the conclusion follows. 
%$E^1\subset Y^{\r\q\n+r}\SUBset E^0$.
\end{pf}  

 
Given two critical puzzle pieces $E^1\subset \inter E^0$, we can construct 
the {\it (Modified) Principle Nest} of critical puzzle pieces 
$$
  E^0\SUPset E^1\SUPset E^1\SUPset\dots\SUPset E^{\chi-1}\SUPset E^\chi 
$$
as described in  in \cite{KL-high}.                        
It comes together with quadratic-like maps $g_n: E^n\ra E^{n-1}$. 

If the map $f$ is renormalizable then  the Principle Nest terminates at some level $\chi$. 
In this case, the last quadratic-like map $g_\chi: E^\chi\ra E^{\chi-1}$  has connected Julia set
that coincides with  the Julia set of the  renormalization $R_\sbb f$, where
$\bbe$ is the $f$-orbit of the non-dividing fixed point $\beta$ of $g_\chi$.  
The renormalization level $\chi$ is also called the {\it height} of the nest. 


% Let $S_i$, $i=i,\dots, l$, be the components of $K(f)\sm \inter Y^0$. 



\subsection{Stars}\label{stars}

% Let $\r\q\n= \k$. 

Given a vertex $v$ of some puzzle piece of depts $n$,
let $S^n(v)$ stand for the union of the puzzle pieces of depth $n$ attached to $v$
(the ``star'' of $v$).
Given a finite set ${\bf v}=\{v_j\}$ of vertices $v_j$, we let
$$
       S^n({\bf v})=\bigcup_j S^n(v_j). 
$$

  Let us begin with an obvious observation that follows from Lemma~\ref{Subset}:

\begin{lem}\label{compactly conained}
  If a puzzle piece $Y^n(z)$ is not contained in $S^n(\balpha)$
then $Y^n(z)\Subset Y^0(z)$.
\end{lem}


\begin{lem}\label{single point}
  For $\bl=\k + 1+2\r$, the stars $S^{\bl}(\alpha_j)$ do not overlap and do not contain the critical point. 
\end{lem}

\begin{pf}
Let us consider the  curves $\CC$ and $f^{-\bt m} \CC'$, $m=1,\dots, \s-1$,  from Lemma \ref{separation}. 
They separate $\alpha'$ from all points of $\balpha\cup \balpha'\sm \{\alpha\}$. 
Furthermore, since $f^{\k}(0)$ is separated from $0$ by $\CC'$,
there is a lift $\Gamma$ of $\CC'$ under $f^{\k}$ that separates $\alpha'$ from $0$ 
and hence from $\alpha$.  It follows that the curves $\Gamma$ and $f^{-\bt m} \CC'$, $m=1,\dots, \s-1$,
separate $\alpha'$ from all points of $\balpha\cup \balpha'$. 
Since the maximal depth of these curves is  $\bl=\k +1$ (which is the depth of $\Gamma$),
the star $S^{\k+1}(\alpha')$ does not overlap with the interior of the stars $S^{\k+1}(a)$
for all  other $a\in \balpha\cup\balpha'$.

By symmetry, the same is true for the star $S^{\k+1}(\alpha)$.
Since these stars  do not contain $0$,
the pullback of $S^{\k+1+\r}(\alpha)$  under $f^\r$ (along $\balpha$) is compactly contained in 
its interior, $\inter S^{\k+1+\r}(\alpha)$.
It follows that $S^{\k+1+\r}(\alpha)$ does overlap with the stars
$S^{\k+1}(a)$ for all  other $a\in \balpha\cup\balpha'$.

Pulling this star once more around $\balpha$, 
we obtain a disjoint family of stars.
Hence all the stars $S^{\k+1+2\r}(a)$,  $a\in \balpha\cup\balpha'$,
are pairwise disjoint. 
\end{pf}

% \bignote{Add to the next section: $S^\bl(\alpha) \Subset S^1(\alpha)$.}


\subsection{Geometric puzzle pieces}

In what follows  we will deal with more general puzzle pieces.

Given  a puzzle piece $Y$, of depth $m$, 
let $Y[l]$ stand for a Jordan disk bounded by the same external rays
as $Y$ and arcs of equipotentials of level $l$ (so $Y[m]=Y$).
Such a  disk will be called a puzzle piece of {\it bidepth} $(m,l)$.
  
A {\it geometric puzzle piece} of bidepth $(m,l)$ is a closed Jordan domain which is 
the union of several puzzle pieces of the same bidepth. 
As for ordinary pieces, a pullback of a geometric puzzle piece of bidepth $(m,l)$ 
under some iterate $f^k$ is a geometric puzzle piece of bidepth $(m+k, l+k)$.
Note also that if $P$ and $P'$ are geometric puzzle pieces with%
\footnote{the inequality between bidepths is understood componentwise}
$\bidepth P \geq \bidepth P'$ 
and $K_P\subset K_{P'}$ then $P\subset P'$.

The family of geometric puzzle pieces of bidepth $(m,l)$ will be called $\YY^m[l]$.
%Given a geometric puzzle piece $Y\in \YY^m[l]$, we let $Y[k]$ be the puzzle piece
%bounded by the same external rays as $Y$ truncated by the equipotential of level $k$. 
%(In particular, $Y[l]=Y$.)

Stars give examples of geometric puzzle pieces. 
Note that pullback of a star is a geometric puzzle piece as well 
but it is a star only if the pullback is univalent.


\comm{*******************8
Any puzzle piece $Y\in \YY^m[l]$ admits the following combinatorial representation.
Let $\theta_i$ be the cyclically ordered  angles of the external rays $\RR_i$ that bound $Y$. 
Let us consider the straight rays $R_i$ in $\C\sm \D$ of angles $\theta_i$
truncated by the circle $\T_r$ of radius $r=1/2^l$. 
If two consecutive rays, $\RR_i$ and $\RR_{i+1}$, land at the same vertex of $Y$,
let is connect $R_i$ to $R_{i+1}$ with a hyperbolic geodesic in $\D$. 
Otherwise $\RR_i$ and $\RR_{i+1}$ are connected with an equipotential arc. 
Then let us connect $R_i$ to $R_{i+1}$ with the appropriate arc of $\T_r$.
We obtain a Jordan curve that bounds the {\it combinatorial model} $M_Y$ of $Y$.

 The arcs $\om_i$ of  $\T\cap M_Y$ correspond to  the 
``external arcs'' of the Julia piece $K_Y$.  
% (which are the different approaches to $K_Y$ from $\infty$ in $\C\sm \cup \bar \RR_i$).
 They have length $2\pi \la$, where $\la$ is called the {\it combinatorial length}
of the corresponding external arc of $K_Y$. 
For a geometric puzzle piece $Y$ of depth $m$,
the combinatorial length of its external arcs is at least $c(\PP)2^{-m}$,
where $c(\PP)$ is the minimum of the combinatorial lengths of the external arcs of the
puzzle pieces of depth 0. 
******************************************}

\msk
% Let $Z=S^1(\alpha)$. 
%Let $L= S^1(\alpha)\sm \inter Y^1$ and $R=-L = S^1(\alpha')\sm \inter Y^1 $. 

\begin{lem}\label{qn-pullbacks}
  Given a point $z\in \inter S^1(\alpha)\cup \inter S^1(\alpha')$ such that $f^{\k} z \in\inter S^1(\alpha)$,
let $P = f^{-\k}(S^1(\alpha))|z$.  Then $P\subset S^1(\alpha)$ or $P\subset S^1(\alpha')$.
\end{lem}

\begin{pf}
  Notice that $\alpha$ and $\alpha'$ are the only points of $\balpha\cup \balpha'$ 
contained in $\inter S^1(\alpha)\cup \inter S^1(\alpha')$.
Since $f^{\k}$ maps $\tl\balpha=\balpha\cup \balpha'\sm \{\alpha, \alpha'\}$
to $\balpha\sm \{\alpha\}$,
no point of $\tl\balpha$ is contained in $\inter P$.
Hence $P$ is contained in $ S^1(\alpha)\cup  S^1(\alpha')$.
%Hence $K_P$ is contained in the closure of some component of $K(f)\sm \tl\balpha$. 
%Since $z\in L$, it must be the closure of the component containing $\alpha$, 
%i.e.,  $K_L\cup K_{Y^1}\cup K_R$.  Hence 
%$$P \subset L\cup Y^1\cup R.$$ 

But by construction of $Y^{\k+1}$, the interior of $f^\k (Y^{\k+1})$ does not overlap with $S^1(\alpha)$.
Hence $\inter Y^{\k+1}$ does not overlap with $P$. 
But $S^1(\alpha)\cup S^1(\alpha')\sm  \inter Y^{\k +1}$ 
consists of two components, one inside $S^1(\alpha)$ and the other inside $S^1(\alpha')$. 
Since $P$ is connected, it is contained in one of them.
\end{pf}

\comm{*********************

Let us now consider the following geometric puzzle pieces:
\begin{itemize}
\item puzzle piece $Z^0=  - Y^0$;
\item  puzzle piece $R$  
   bounded by the arcs $f^{-\bt m} (\CC')$, $m=0,1,\dots, \s-1$, from Lemma \ref{separation} 
   and the equipotential of level one; 
\item the symmetric puzzle piece $L=-R$.
\end{itemize}


\begin{lem}\label{qn-pullbacks}
  Let $z\not\in S^{3\r}(\balpha)$,  $g^{\q\n} z \in\inter Z^0$ and let $P = g^{-\q\n}(Z^0)|z$. 
Then $P\Subset \inter Y^0$ or  $P\Subset \inter Z^0$.
\end{lem}

\begin{pf}
  $P$ is a geometric puzzle piece of bidepth $(\r\n\q+1, \r\n\q) $.
 But the puzzle piece $Y^{\r\n\q+1}$ is mapped under $g^{\n\q}$ into  $R$, 
  and $R \cap Z^0=\emptyset$.
It follows that $P \cap \inter Y^{\r\n\q+1}=\emptyset$. 
But $K(f)\sm \inter Y^{\r\q\n+1}$ consists of two $0$-symmetric 
connected components $X_L\supset L\cap K(f)$ and $X_R\supset R\cap K(f)$. 
We conclude that  $K_P$ is contained in one of these components, and hence it is contained in one of the
sets $K_{Z^0}$ or $K_{Y^0}$. As
 $$
    \bidepth P\geq (1,0)=\bidepth Z^0\geq (0,0)=\bidepth Y^0,
$$ 
 $P$ is contained in one of the
puzzle pieces $Z^0$ or $Y^0$. 
\end{pf}
*******************************}
 
\comm{*************
\begin{lem}\label{returns to L}
  If  $z \in S^{3\r} (\balpha)$ then $f^k z\in L$
for some $k\leq 2\r$. 
\end{lem}

\begin{pf}
   If $z\in S^{3\r}(\alpha_j)$, then 
$
  \zeta=f^{\r-j}z\in S^{2\r}(\alpha)\subset L\cup Y^\r .
$
But if $\zeta\in Y^\r$ then $f^\r \zeta\in L$.    
\end{pf}
**************************}


% \section {Geometric separation between $\ba$ and $\ba'$}








